\documentclass[tikz,border=10pt]{standalone}
\usetikzlibrary{positioning, arrows.meta, calc, shapes.geometric, babel}
\usepackage[utf8]{inputenc}
\usepackage[T1]{fontenc}
\usepackage{lmodern}
\usepackage[spanish,es-tabla]{babel}
\usepackage{amsmath, amssymb, bm}
\spanishdecimal{.}

\begin{document}
\begin{tikzpicture}[
  node distance=0.8cm and 1.0cm,
  align=center,
  font=\sffamily\small,
  % Estilos de bloques
  box/.style={draw, rectangle, rounded corners=3pt, minimum width=3.5cm, minimum height=1.1cm, thick, fill=white},
  decision/.style={draw, diamond, aspect=1.8, thick, fill=white, inner sep=2pt},
  discardbox/.style={box, fill=red!5, draw=red!70, dashed},
  successbox/.style={box, fill=green!5, draw=green!70},
  % Estilos de flechas
  line/.style={draw, -{Stealth[scale=1.2]}, thick},
  dashline/.style={draw, -{Stealth[scale=1.2]}, thick, dashed}
]

  % ==========================================
  % BLOQUE INICIAL: CONFIGURACIÓN
  % ==========================================
  \node[box, fill=gray!5] (inicio) {\textbf{Configuración de Escenario}\\Definir trayectoria nominal\\y ganancias base};

  % ==========================================
  % FASE 1: DIAGNÓSTICO
  % ==========================================
  \node[box, below=of inicio] (diagnostico) {
    \textbf{Fase 1: Diagnóstico Inicial}\\
    Evaluar ganancias base en bucle cerrado
  };
  \draw[line] (inicio) -- (diagnostico);

  \node[decision, below=0.8cm of diagnostico] (dec_base) {{\textquestiondown}Supera umbral\\de RMSE?};
  \draw[line] (diagnostico) -- (dec_base);

  % ==========================================
  % FASE 2: BÚSQUEDA GLOBAL
  % ==========================================
  \node[box, below=0.9cm of dec_base] (global) {
    \textbf{Fase 2: Búsqueda Global}\\
    Generar 32 candidatos aleatorios\\
    mediante muestreo log-uniforme
  };
  % Conexión si NO supera el umbral (requiere tuneo)
  \draw[line] (dec_base) -- node[right, font=\sffamily\bfseries\color{red!70!black}] {NO (sintonización)} (global);

  % ==========================================
  % FASE 3: FILTROS LÍMITE
  % ==========================================
  \node[box, below=of global] (filtros) {
    \textbf{Fase 3: Filtros Límite}\\
    Evaluar en bucle cerrado
  };
  \draw[line] (global) -- (filtros);

  % Criterios de filtro
  \node[decision, below=0.8cm of filtros] (dec_filtros) {{\textquestiondown}Estable y seguro?};
  \draw[line] (filtros) -- (dec_filtros);

  % Descartes y Aceptación
  \node[discardbox, left=1.2cm of dec_filtros] (descarte) {
    \textbf{Descarte} [$\otimes$]\\
    - Caídas ($z < z_{\min}$)\\
    - Exceso inclinación\\
    - Inestabilidad/NaNs\\
    - Saturación $> 2.0$ s
  };
  \draw[line] (dec_filtros.west) -- node[above, font=\sffamily\bfseries\color{red!70!black}] {NO} (descarte.east);

  \node[successbox, below=0.9cm of dec_filtros] (supervivientes) {
    \textbf{Supervivientes} [$\checkmark$]\\
    Selección de hasta 16 mejores\\
    candidatos iniciales
  };
  \draw[line] (dec_filtros) -- node[right, font=\sffamily\bfseries\color{green!60!black}] {SÍ} (supervivientes);

  % ==========================================
  % FASE 4: REFINAMIENTO LOCAL
  % ==========================================
  \node[box, below=of  supervivientes] (local) {
    \textbf{Fase 4: Refinamiento Local}\\
    Tipo simplex / Nelder--Mead\\
    para minimizar RMSE
  };
  \draw[line] (supervivientes) -- (local);

  % ==========================================
  % FASE 5: SELECCIÓN MULTICRITERIO
  % ==========================================
  \node[box, below=of local] (seleccion) {
    \textbf{Fase 5: Selección y Desempate}\\
    1. Minimizar RMSE pos/vel\\
    2. Penalizar esfuerzo de actuación\\
    3. Penalizar desviación de ganancias base
  };
  \draw[line] (local) -- (seleccion);

  % ==========================================
  % RESULTADO FINAL: YAML CONGELADO
  % ==========================================
  \node[box, fill=gray!10, below=1.0cm of  seleccion, minimum width=4.5cm] (yaml) {
    \textbf{Resultado Final}\\
    Ganancias ganadoras fijadas
  };
  \draw[line] (seleccion) -- (yaml);

  % Conexión directa desde Fase 1 (SÍ supera el umbral)
  \draw[line] (dec_base.east) -- node[above, xshift=0.3cm, font=\sffamily\bfseries\color{green!60!black}] {SÍ} ++(2.8,0) |- (yaml.east);

\end{tikzpicture}
\end{document}
